%========================================================
% CHAPTER 2 — Modeling Strategy
%========================================================

\chapter{Modeling Strategy}

%========================================================
% SECTION 2.1 — Progressive Modeling Philosophy
%========================================================

\section{Progressive Modeling Philosophy}

%--------------------------------------------------------
% Subsection 2.1.1 — Progressive Modeling Approach
%--------------------------------------------------------

\subsection{Progressive Modeling Approach}

The modeling approach follows a progressive refinement strategy, in which the system is described through successive levels of approximation of increasing physical complexity.

Each stage introduces additional physical effects while remaining consistent with the previous ones. This ensures analytical transparency and preserves continuity between simplified and refined models.

\medskip

From a theoretical perspective, this approach isolates the contribution of individual mechanisms. From a computational perspective, it enables a layered simulator structure, where each refinement corresponds to a well-defined extension of the existing codebase.

%--------------------------------------------------------
% Subsection 2.1.2 — Hierarchy of Modeling Levels
%--------------------------------------------------------

\subsection{Hierarchy of Modeling Levels}

The modeling workflow is structured as a sequence of progressively refined descriptions:

\[
\begin{array}{c}
\text{Ideal Bell State} \\
\downarrow \\
\text{Local Phase Model (Z-rotations)} \\
\downarrow \\
\text{General Local Unitary Transformations} \\
\downarrow \\
\text{Mixed-State Description (Density Operators)} \\
\downarrow \\
\text{Noisy Quantum Channels} \\
\downarrow \\
\text{Effective Fiber Models}
\end{array}
\]

\medskip

More explicitly:

\begin{itemize}

\item \textbf{Level 1: Minimal phase model}  
Pure, maximally entangled state with fiber action reduced to a relative phase. This level captures the analytical dependence of correlation observables on a single parameter.

\item \textbf{Level 2: General unitary dynamics}  
Arbitrary polarization rotations described by local operators \( U_A \otimes U_B \).

\item \textbf{Level 3: Mixed-state description}  
Density operators used to represent statistical mixtures and loss of coherence.

\item \textbf{Level 4: Noisy quantum channels}  
Effective completely positive trace-preserving (CPTP) maps modeling depolarization and decoherence.

\item \textbf{Level 5: Effective fiber models}  
Physically motivated descriptions, such as concatenations of random unitary segments or stochastic polarization evolution.

\end{itemize}

\medskip

This hierarchy defines the conceptual roadmap of the present work.

The initial levels remain analytically tractable and allow direct validation, while the later levels progressively increase physical realism and connect to experimentally relevant propagation effects.

%========================================================
% SECTION 2.2 — Baseline Quantum Model
%========================================================

\section{Baseline Quantum Model}

%--------------------------------------------------------
% Subsection 2.2.1 — Pure-State Description
%--------------------------------------------------------

\subsection{Pure-State Description}

At the first stage, the bipartite system is described as a pure state in \( \mathbb{C}^2 \otimes \mathbb{C}^2 \), and fiber propagation is modeled through local unitary transformations.

Particular emphasis is placed on relative phase shifts induced by the fibers.

\medskip

The main objective is to analyze how these phase shifts affect measurable two-photon correlations.

The relevant observables are operators of the form
\[
\sigma_i \otimes \sigma_j,
\qquad
i,j \in \{x,y,z\},
\]
as well as operators corresponding to arbitrary measurement directions on the Bloch sphere.

\medskip

These operators play a central role because they describe joint local measurements performed independently on the two subsystems.

More precisely, the operator \( \sigma_i \) acts on subsystem \( A \), while \( \sigma_j \) acts on subsystem \( B \). Their tensor product therefore represents a correlation measurement between the outcomes obtained along two specified measurement directions.

\medskip

The corresponding expectation values
\[
\langle \sigma_i \otimes \sigma_j \rangle
\]
quantify the statistical correlations between local measurements performed on the two photons.

These quantities contain direct information about the nonlocal structure of the bipartite quantum state and therefore constitute one of the fundamental signatures of entanglement.

\medskip

For polarization-entangled states, these observables are particularly important because they are directly connected to experimentally accessible coincidence measurements.

In the computational basis, the Pauli operators correspond to measurements along specific polarization bases:
\begin{itemize}
\item \( \sigma_z \): horizontal/vertical basis,
\item \( \sigma_x \): diagonal/antidiagonal basis,
\item \( \sigma_y \): circular polarization basis.
\end{itemize}

Consequently, expectation values of the form
\[
\langle \sigma_i \otimes \sigma_j \rangle
\]
can be interpreted as experimentally measurable polarization correlations between the two photons.

\medskip

Beyond their direct physical interpretation, these correlators also provide the building blocks for more advanced mathematical descriptions introduced later in this work.

In particular:
\begin{itemize}
\item the diagonal correlations \( \langle \sigma_i \otimes \sigma_i \rangle \) characterize the persistence of polarization correlations under fiber propagation,
\item the complete set of correlators defines the correlation tensor representation of the bipartite state,
\item combinations of these expectation values enter directly into Bell inequalities and CHSH nonlocality tests.
\end{itemize}

\medskip

This baseline model therefore constitutes the simplest regime in which entanglement is preserved and can be probed exclusively through correlation observables.

It provides a controlled setting for deriving analytical expressions and validating the numerical implementation.

\medskip

This regime establishes the minimal framework required to connect fiber-induced phase evolution with measurable two-photon correlations, while simultaneously introducing the correlation structures that will later be generalized to mixed states, arbitrary measurement axes, and non-unitary evolution.

%--------------------------------------------------------
% Subsection 2.2.2 — Extension to Mixed-State Formalism
%--------------------------------------------------------

\subsection{Extension to Mixed-State Formalism}

At the next level, the model incorporates statistical effects and decoherence.

In this regime, the pure-state description is no longer sufficient and must be replaced by a density-operator formalism.

\medskip

This extension enables the description of:
\begin{itemize}
\item reduced states of individual photons,
\item loss of coherence due to unresolved degrees of freedom,
\item effective non-unitary dynamics arising from environmental interactions or statistical averaging.
\end{itemize}

\medskip

Within this framework, physically relevant quantities such as purity, fidelity, and concurrence can be defined and computed, enabling a quantitative characterization of mixedness, coherence, and entanglement (see Chapter~4).

\medskip

A central mathematical operation in the density-operator formalism is the partial trace, which allows the description of individual subsystems starting from the global bipartite state.

For a composite system described by the density operator \( \rho_{AB} \), the reduced state of subsystem \( A \) is obtained by tracing over subsystem \( B \):

\begin{equation}
\rho_A
=
\mathrm{Tr}_B(\rho_{AB})
\label{eq:partial_trace_A}
\end{equation}

\medskip

Analogously, the reduced state of subsystem \( B \) is

\begin{equation}
\rho_B
=
\mathrm{Tr}_A(\rho_{AB}).
\label{eq:partial_trace_B}
\end{equation}

\medskip

The reduced density operators contain all information accessible through local measurements performed on the corresponding subsystems.

\medskip

For maximally entangled Bell states such as \( |\Psi^+\rangle \), the reduced states are maximally mixed:

\begin{equation}
\rho_A
=
\rho_B
=
\frac{1}{2} I.
\label{eq:maximally_mixed_reduced_states}
\end{equation}

This result constitutes one of the fundamental signatures of bipartite entanglement.

Although the global state is pure, each subsystem individually admits a mixed-state description, implying that no local measurement can fully reconstruct the global quantum state.

The information characterizing the entangled system is therefore encoded nonlocally in the correlations between the two subsystems.

\medskip

This connection between entanglement, reduced-state mixedness, and nonlocal correlations plays a central role throughout the present work, both in the theoretical analysis and in the interpretation of the numerical simulations developed in later chapters.
\medskip

Expectation values in the density-operator formalism are computed as

\begin{equation}
\langle O \rangle = \mathrm{Tr}(\rho O)
\label{eq:density_expectation}
\end{equation}

which generalizes the pure-state expression \( \langle \psi | O | \psi \rangle \).

\medskip

The density-operator formalism therefore provides the natural framework for describing ensemble averaging, decoherence, and effective non-unitary dynamics arising during fiber propagation.

%========================================================
% SECTION 2.3 — Computational Architecture and Modeling Roadmap
%========================================================

\section{Computational Architecture and Modeling Roadmap}

%--------------------------------------------------------
% Subsection 2.3.1 — Simulator Design and Architecture
%--------------------------------------------------------

\subsection{Simulator Design and Architecture}

The simulator is designed as a modular implementation of the theoretical model.

\medskip

The initial modules include:
\begin{itemize}
\item Bell-state generation,
\item local unitary propagation,
\item correlation evaluation.
\end{itemize}

\medskip

More advanced modules introduce:
\begin{itemize}
\item quantum channels,
\item reduced density operators,
\item entanglement measures,
\item experimentally relevant observables such as Bell parameters and Stokes-like quantities.
\end{itemize}

\medskip

The theoretical structure is directly reflected in the computational architecture, ensuring consistency, traceability, and extensibility.

\medskip

The implementation distinguishes clearly between reusable computational kernels and higher-level experiment scripts, allowing the simulator to evolve into a structured and reusable research tool.

%--------------------------------------------------------
% Subsection 2.3.2 — Implementation Strategy
%--------------------------------------------------------

\subsection{Implementation Strategy}

The simulator is designed as a layered computational framework mirroring the progressive hierarchy of the theoretical model.

\medskip

At the initial level, the implementation operates on pure two-qubit states represented as complex vectors in \( \mathbb{C}^4 \).

The corresponding modules include:
\begin{itemize}
\item Bell-state generation,
\item local unitary evolution,
\item expectation-value evaluation,
\item correlation computation.
\end{itemize}

\medskip

The extension to mixed-state dynamics introduces:
\begin{itemize}
\item density-matrix representations,
\item partial-trace operations,
\item observable evaluation in the density formalism,
\item state metrics such as purity, fidelity, and concurrence.
\end{itemize}

\medskip

More advanced layers incorporate:
\begin{itemize}
\item quantum channels,
\item ensemble-based evolution,
\item stochastic propagation models,
\item experimentally motivated observables.
\end{itemize}

\medskip

At the current stage, this structure is reflected in the repository organization:
\begin{itemize}
\item \texttt{Core}: state vectors, operators, expectation values, and correlation routines;
\item \texttt{Models}: physical propagation models;
\item \texttt{Utils}: validation routines and auxiliary tools;
\item \texttt{Experiments}: scripts implementing complete simulation studies.
\end{itemize}

\medskip

This modular architecture preserves analytical traceability while enabling incremental extensions of the physical model without structural modifications of previously validated components.